
\documentclass[tikz,border=2pt]{standalone}
\usepackage[utf8]{inputenc}
\usepackage[T1]{fontenc}
\usepackage[french]{babel}
\usepackage{amsmath,amssymb}
\usepackage[table]{xcolor}
\usepackage{siunitx}
\usepackage[version=4]{mhchem}
\usepackage{tabularx,booktabs,array}
\usepackage{tikz}
\usetikzlibrary{arrows.meta,positioning,calc,fit,decorations.pathmorphing,patterns,shapes.geometric}
\usepackage{pgfplots}
\pgfplotsset{compat=1.18}
\definecolor{primary}{RGB}{14,116,144}
\definecolor{accent}{RGB}{16,185,129}
\definecolor{dark}{RGB}{15,23,42}
\definecolor{bglight}{RGB}{241,245,249}
\definecolor{borderlight}{RGB}{203,213,225}
\definecolor{examplepink}{RGB}{236,72,153}
\definecolor{remarkblue}{RGB}{14,165,233}
\definecolor{notepurple}{RGB}{99,102,241}
\definecolor{synthgreen}{RGB}{5,150,105}
\definecolor{synthorange}{RGB}{249,115,22}
\definecolor{synthviolet}{RGB}{79,70,229}
\definecolor{primary}{RGB}{14,116,144}
\definecolor{accent}{RGB}{16,185,129}
\definecolor{dark}{RGB}{15,23,42}
\definecolor{bglight}{RGB}{241,245,249}
\definecolor{examplepink}{RGB}{236,72,153}
\definecolor{remarkblue}{RGB}{14,165,233}
\definecolor{notepurple}{RGB}{99,102,241}
\definecolor{synthgreen}{RGB}{5,150,105}
\definecolor{synthorange}{RGB}{249,115,22}
\definecolor{synthviolet}{RGB}{79,70,229}
\begin{document}
\begin{tikzpicture}[scale=0.85]
 % Paroi avec fente
 \fill[gray!60] (-4,1.8) rectangle (-0.4,2.2);
 \fill[gray!60] (0.4,1.8) rectangle (4,2.2);
 
 % Etiquette de la fente
 \draw[<->, thick, primary] (-0.4,2.5) -- (0.4,2.5);
 \node[above, primary, font=\small\bfseries] at (0,2.5) {$a$};
 
 % Ondes incidentes (lignes droites)
\foreach \y in {-1.5,-1.0,-0.5,0,0.5,1.0,1.5} {
\draw[blue!60, thick] (-4,\y) -- (-0.5,\y);
}
 % Ondes incidentes avant la paroi
 \foreach \y in {-1.5,-1.0,-0.5,0,0.5,1.0,1.5} {
 \draw[blue!50, thick] (-4,\y) -- (-0.45,\y);
 }
 
 % Fleche direction incidente
\draw[-{Stealth[length=3mm]}, very thick, dark] (-3.5,-2.5) -- (-1.2,-2.5);
\node[below, dark, font=\small] at (-2.35,-2.5) {Direction incidente};
 
 % Ondes diffractees (arcs de cercle)
 \foreach \r in {1.0,1.5,2.0,2.5,3.0,3.5} {
 \draw[red!60, thick] (0,2) ++(90-60:\r) arc (90-60:90+60:\r);
 }
 
 % Fleches de directions de propagation diffractees
\draw[-{Stealth[length=2.5mm]}, very thick, red!70!black] (0,2) -- (1.8,5.2);
\draw[-{Stealth[length=2.5mm]}, very thick, red!70!black] (0,2) -- (0,5.5);
\draw[-{Stealth[length=2.5mm]}, very thick, red!70!black] (0,2) -- (-1.8,5.2);
\draw[dashed, dark] (0,2) -- (0,5.0);
\draw[red!70!black, thick] (0,3.0) arc[start angle=90,end angle=60,radius=1.0];
\node[red!70!black, font=\small\bfseries] at (0.45,3.35) {$\theta$};
 
 \node[right, red!70!black, font=\small, text width=3cm] at (2.2,4.5) {\'{E}talement des directions (diffraction)};
 
 % Labels
 \node[below, font=\small\itshape, gray!80!black] at (0,1.6) {Fente de largeur $a$};
 \node[left, font=\small, blue!70!black] at (-4,0) {Ondes incidentes};
\end{tikzpicture}
\end{document}
